% !PRINTED 455
\weberpage{455}

%<<CONT
\gesperrt{able, those of the first kind being designated as indecomposable functions}. If, for example, $\Omega$ is the field of rational numbers, then $x^2 - y^2 = (x - y)(x + y)$ is reducible in $\Omega$. The function $x^2 - 2y^2$ is indeed decomposable into the factors $x + y\sqrt{2}, x - y\sqrt{2}$, but these factors are not functions in $\Omega$. The function $x^2 + y^2 + 1$, finally, is not decomposable at all. A function irreducible in a field $\Omega$, but decomposable, becomes reducible in an extended field.

The fundamental theorems on irreducible functions we came to know in \S.\,51. In those developments, it is true, the field of all numbers was at first assumed as the domain of rationality, in accordance with the standpoint there taken, in which field the distinction between indecomposable and irreducible functions disappears. But only the reducibility or irreducibility of a function is ever made use of, and the particular nature of the field $\Omega$ does not come into consideration. Thus there is contained in it the proof of the theorem:

\gesperrt{IV. A reducible function in $\Omega$ can be decomposed into irreducible factors in only one way, if rational functions which differ only by a constant factor are regarded as not distinct.}

\begin{center}
§. 142.\\
\gesperrt{Algebraic Fields.}
\end{center}

If $\Omega$ is an arbitrary field, and $F(x)$ a function in $\Omega$, then the equation
\weq{1}{F(x) = 0}
is called an \gesperrt{equation in} $\Omega$. This equation is called reducible or irreducible according as the function $F(x)$ is reducible or irreducible. By the adjunction of a root $\alpha$ of such an equation to $\Omega$ there arises (if $\alpha$ does not itself already belong to $\Omega$) a new field $\Omega'$, which we shall call an \gesperrt{algebraic field} ``\gesperrt{over}'' $\Omega$, or also, when doubt as to the meaning is excluded, briefly, an \gesperrt{algebraic field}. For such a field we use the sign
\[ \Omega(\alpha). \]
